\documentclass[11pt]{article}

\usepackage[T1]{fontenc}
\usepackage[utf8]{inputenc}
\usepackage{lmodern}
\usepackage[margin=1in]{geometry}
\usepackage{microtype}
\usepackage{amsmath,amssymb,amsthm,mathtools}
\usepackage{booktabs}
\usepackage[numbers,sort&compress]{natbib}
\usepackage{hyperref}
\hypersetup{
  colorlinks=true,
  linkcolor=black,
  citecolor=black,
  urlcolor=blue
}

\newcommand{\Nat}{\mathbb{N}}
\newcommand{\ZZ}{\mathbb{Z}}
\newcommand{\RR}{\mathbb{R}}
\newcommand{\gn}[1]{\ulcorner #1 \urcorner}
\DeclareMathOperator{\Prov}{Prov}
\DeclareMathOperator{\Con}{Con}
\DeclareMathOperator{\prf}{Proof}
\newcommand{\PA}{\mathsf{PA}}
\newcommand{\ZFC}{\mathsf{ZFC}}
\newcommand{\GL}{\mathsf{GL}}
\newcommand{\GLS}{\mathsf{GLS}}

\newtheorem{theorem}{Theorem}
\newtheorem{proposition}[theorem]{Proposition}
\newtheorem{lemma}[theorem]{Lemma}
\newtheorem{corollary}[theorem]{Corollary}
\theoremstyle{definition}
\newtheorem{definition}[theorem]{Definition}
\theoremstyle{remark}
\newtheorem{remark}[theorem]{Remark}

\title{Self-Certifying Return%
\thanks{Packaging note for the Eternal Return table.
The recurrence theorem is Poincar\'e's, the incompleteness theorems are G\"odel's,
the fixed-point theorem for provability is L\"ob's, and the modal completeness theorem is Solovay's.
The note only keeps their quantifiers from being exchanged.
It is the formal companion to the table review in the repository root.}}
\author{Eternal Return table}
\date{28 September 2026}

\begin{document}
\maketitle

\begin{abstract}
This note separates Poincar\'e recurrence from self-certifying readings of ``this life returns.''
In a measure-preserving system of finite measure, almost every point of a positive-measure set returns to that set infinitely often; when the space is separable and metric, almost every point returns to every neighborhood of itself (\(\varepsilon\)-Return).
Exact identity of a continuous state is a different statement, and it is typically not a predicate of a recursively axiomatized arithmetic \(S\): \(S\) can express equality of codes and coded neighborhood membership.
Let \(\Box\) be provability in a standard proof predicate for a consistent, recursively axiomatized extension \(S\) of Peano arithmetic.
The logic of \(\Box\) is G\"odel--L\"ob logic \(\GL\).
The scheme \(\Box A \to A\) is not a theorem of \(\GL\); L\"ob's theorem replaces it, and the L\"ob hypothesis is not the hypothesis of Poincar\'e's theorem.
A configuration that encodes \(S\), or encodes a proof of its own return, is a self-referential certificate.
The Henkin sentence for \(S\) is provable in \(S\) and does not prove \(\Con_S\).
A fixed point that conjoins a recurrence formula with \(\Con_S\) is unprovable if \(S\) is consistent.
These are limits on self-certification inside \(S\).
They do not say that recurrent sets are empty.
Separately, recurrence queries for configurations that encode a universal computation are undecidable.
Undecidability of the query does not decide whether a given orbit returns.
Nothing here equates \(\varepsilon\)-Return with the demand to will an identical life, and incompleteness of \(S\) is neither a refutation of the recurrence theorem nor an argument for it.
\end{abstract}

\section{Introduction}

Three different claims travel under the word ``return.''
This note gives each one a sentence, a system, and a hypothesis list.

The first is dynamical.
A measure-preserving transformation of a finite measure space returns almost every point of a positive-measure set to that set, infinitely often \citep{poincare1890,walters1982}.
In a separable metric space the same argument yields \emph{\(\varepsilon\)-Return}: almost sure return into every neighborhood of the point.
The second is linguistic.
A formal arithmetic \(S\) does not contain an arbitrary point of a continuum as a value of a variable.
It contains codes.
The predicates it can form of those codes are code-identity and coded neighborhood membership.
The third is self-certifying.
One asks the configuration to contain the system \(S\), or a proof in \(S\), and to count that included proof as establishing the return of the configuration that contains it.
The tools for that reading are the second incompleteness theorem \citep{godel1931}, L\"ob's theorem \citep{lob1955}, and the provability logic \(\GL\) \citep{solovay1976,boolos1993}.
Throughout, \(\Box\) means provability in one named system \(S\), not an unspecified necessity and not a knower.

The note proves no new incompleteness theorem and no new recurrence theorem.
Section~\ref{sec:setup} records \(\varepsilon\)-Return with its hypotheses and the predicate boundary.
Section~\ref{sec:modal} packages \(\Box\) as \(\GL\).
Section~\ref{sec:epistemic} says why an \(\mathsf{S5}\) knowledge operator is a different symbol.
Section~\ref{sec:barrier} applies the second theorem and L\"ob's theorem to two named fixed points.
Section~\ref{sec:undecidable} records a halting reduction for the query ``does \(C\) recur?''
Section~\ref{sec:not} lists the readings that these results do not support.
G\"odel's rotating cosmological models \citep{godel1949} are a solution of Einstein's equations with closed timelike curves.
They are not an argument about \(\Prov_S\), and they are not used below.

\section{Setup: dynamics and the formal language}
\label{sec:setup}

Two systems are in play, and they are not the same object.
The dynamical statement lives in analysis.
The provability statement lives in arithmetic.

\begin{definition}[Foundations]
\label{def:systems}
\(\ZFC\) is Zermelo--Fraenkel set theory with choice, used only as a sufficient ambient theory in which Lebesgue measure and the argument of Theorem~\ref{thm:poincare} can be formalized.
\(S\) is a consistent, recursively axiomatized extension of Peano arithmetic \(\PA\).
The symbol \(\Box\) is never used for provability in \(\ZFC\).
When provability in \(\ZFC\) is intended it is written \(\Prov_{\ZFC}\).
\end{definition}

Consistency of \(S\) is an external assumption about \(S\), in the usual sense that no \(S\)-derivation ends in \(0=1\).
It is not an axiom of \(S\).

\subsection{Measure-preserving dynamics}

\begin{definition}
\label{def:mps}
A \emph{measure-preserving system} is a quadruple \((X,\mathcal{B},\mu,T)\) where \((X,\mathcal{B},\mu)\) is a measure space with \(\mu(X)<\infty\) and \(T\colon X\to X\) is measurable and satisfies \(\mu(T^{-1}A)=\mu(A)\) for every \(A\in\mathcal{B}\).
\end{definition}

Finite measure is the hypothesis that does the work.
Normalizing to a probability does not change the argument.
In the classical physical reading, \(T\) is the time-\(1\) map of a flow that preserves a finite measure on its region (Liouville measure on a finite-volume invariant set is the usual source of such a \(T\)).
The note proves the discrete statement.
A flow inherits integer-time returns from its time-\(1\) map.

\begin{theorem}[Poincar\'e recurrence]
\label{thm:poincare}
Let \((X,\mathcal{B},\mu,T)\) be a measure-preserving system and let \(A\in\mathcal{B}\) satisfy \(\mu(A)>0\).
Then
\[
\mu\bigl(\{x\in A \mid \text{there are infinitely many } n\geq 1 \text{ with } T^{n}x\in A\}\bigr)
= \mu(A).
\]
\end{theorem}

\begin{proof}[Standard argument]
Let \(E=\{x\in A \mid \forall n\geq 1\ (T^{n}x\notin A)\}\).
The sets \(T^{-k}E\) for \(k\geq 0\) are pairwise disjoint.
If \(x\in T^{-i}E\cap T^{-j}E\) for \(0\leq i<j\), then \(T^{i}x\in E\subseteq A\) and \(T^{j-i}(T^{i}x)\in A\), contradicting \(T^{i}x\in E\).
Invariance gives \(\mu(T^{-k}E)=\mu(E)\).
Finite measure gives \(\sum_{k\geq 0}\mu(E)\leq \mu(X)<\infty\), so \(\mu(E)=0\).
Thus almost every point of \(A\) returns to \(A\) at least once.

The same argument applies to each iterate \(T^{N}\), which is measure-preserving.
The set of points of \(A\) that fail to return under every positive iterate of \(T^{N}\) has measure zero, and the set of points that never return at a time \(n\geq N\) is contained in it.
A countable union over \(N\) of null sets is null.
\end{proof}

The theorem is a statement about almost every point of a named measurable set.
It does not say that every point returns, and it does not say that the return time is short.
Both strengthenings need hypotheses this argument does not have.

\begin{definition}[\(\varepsilon\)-Return]
\label{def:epsilon}
Suppose in addition that \(X\) is a separable metric space, \(\mathcal{B}\) contains the Borel \(\sigma\)-algebra, and \(d\) is the metric.
The system has \emph{\(\varepsilon\)-Return} when, for every \(\varepsilon>0\),
\[
\mu\bigl(\{x\in X \mid \text{there are infinitely many } n\geq 1 \text{ with } d(T^{n}x,x)<\varepsilon\}\bigr)
= \mu(X).
\]
\end{definition}

\begin{corollary}
\label{cor:epsilon}
Every measure-preserving system on a separable metric space, with \(\mu\) a finite Borel measure, has \(\varepsilon\)-Return.
\end{corollary}

\begin{proof}
Let \(\{U_i\}_{i\in\Nat}\) be a countable basis of nonempty open sets.
By Theorem~\ref{thm:poincare}, for each \(i\) there is a null set \(N_i\subseteq U_i\) of points that fail to return to \(U_i\) infinitely often.
Let \(N=\bigcup_i N_i\).
If \(x\notin N\) and \(\varepsilon>0\), separability supplies a basis element \(U_i\) with \(x\in U_i\subseteq\{y\mid d(y,x)<\varepsilon\}\).
Infinitely many returns of \(x\) to \(U_i\) are returns to that neighborhood.
The null set \(N\) does not depend on \(\varepsilon\).
\end{proof}

\(\varepsilon\)-Return is neighborhood membership for almost every point.
Exact identity \(T^{n}x=x\) is a different predicate, and the finite-measure hypotheses do not force it.

\begin{proposition}[Exact return can fail everywhere]
\label{prop:rotation}
Let \(X=\RR/\ZZ\), let \(\mu\) be Lebesgue measure, and let \(T(x)=x+\alpha \bmod 1\) with \(\alpha\) irrational.
Then \((X,\mu,T)\) is a measure-preserving system, every point returns to every neighborhood of itself infinitely often, and \(T^{n}x\neq x\) for every \(x\) and every \(n\geq 1\).
\end{proposition}

\begin{proof}
Translation preserves Lebesgue measure, and \(\mu(X)=1\).
If \(T^{n}x=x\), then \(n\alpha\in\ZZ\), so \(\alpha\) is rational.
For the neighborhood claim, write \(\|\,\cdot\,\|\) for distance to the nearest integer, so \(d(T^{n}x,x)=\|n\alpha\|\).
The \(N+1\) points \(0,\{\alpha\},\ldots,\{N\alpha\}\) on the circle cannot be pairwise more than \(1/N\) apart, so two of them lie at distance at most \(1/N\).
Their difference is an integer \(q\) with \(1\leq q\leq N\) and \(\|q\alpha\|\leq 1/N\).
Fix \(\varepsilon>0\) and suppose \(F=\{n\geq 1:\|n\alpha\|<\varepsilon\}\) is finite.
Since \(\alpha\) is irrational, \(\|n\alpha\|>0\) for every \(n\geq 1\).
If \(F\) is empty, any integer \(N>1/\varepsilon\) yields a \(q\) with \(\|q\alpha\|<\varepsilon\), so \(q\in F\), contradicting emptiness.
If \(F\) is nonempty, let \(\delta=\min_{n\in F}\|n\alpha\|>0\) and choose an integer \(N>\max(1/\varepsilon,1/\delta)\).
The same step yields \(q\leq N\) with \(\|q\alpha\|\leq 1/N<\delta\), so \(q\in F\) and \(\|q\alpha\|<\delta\), which contradicts the choice of \(\delta\).
Thus \(F\) is infinite.
Every point returns to every one of its neighborhoods infinitely often, and the system has \(\varepsilon\)-Return with empty exceptional set.
This approximation argument is special to the example.
The measure theorem that applies to every system in Definition~\ref{def:mps} is Theorem~\ref{thm:poincare}, which is almost-everywhere.
\end{proof}

Proposition~\ref{prop:rotation} is the separator between the two dynamical predicates.
The system meets every hypothesis of Corollary~\ref{cor:epsilon}.
The set of exactly periodic points is empty.
A theorem that delivers \(\varepsilon\)-Return has not delivered bit-identity of states.

\subsection{What \texorpdfstring{\(S\)}{S} can say}

Points of \(\RR^{d}\) are not natural numbers.
The variables of \(S\) range over numbers, codes of formulas, and codes of proofs.
A continuous state enters a sentence of \(S\) only after a code has been chosen: a G\"odel number of a definition, an index of a computable Cauchy name, or a code of a basic open set.

\begin{definition}[Code predicates]
\label{def:codes}
Fix a standard primitive recursive G\"odel numbering.
For codes \(c,c'\) and a positive rational \(q\) given by a code, the following are formulas of arithmetic:
\begin{itemize}
\item \(\mathrm{EqCode}(c,c')\), the atomic equality \(c=c'\) of the codes as numbers;
\item \(\mathrm{Near}(c,c',q)\), a formalized statement that the basic open sets (or Cauchy names) numbered by \(c\) and \(c'\) meet within tolerance \(q\).
\end{itemize}
Neighborhood membership, once the neighborhood is coded, is of the first kind of statement.
Identity of the underlying real, when both objects are presented by names, is agreement of those names, which is again a statement about the codes.
Identity of binary expansions is not well-defined on reals: dyadic rationals have two expansions.
\end{definition}

\begin{remark}[Predicate boundary]
\label{rem:boundary}
Exact bit-identity of an uncoded continuous state is not an \(S\)-predicate.
There is no formula \(\varphi(v)\) of the language of \(S\) whose intended extension is a singleton \(\{x\}\) for an arbitrary \(x\in\RR^{d}\), because there is no variable in that language for \(x\).
What \(S\) states, after a coding has been fixed, is \(\mathrm{EqCode}\) or \(\mathrm{Near}\).
Theorem~\ref{thm:poincare}, formalized in \(\ZFC\), is a sentence about measures and preimages.
It is not a scheme of \(S\)-atoms ``\(x\) itself returns,'' one for each real \(x\).
\end{remark}

A symbolic configuration that is already a natural number, or a finite word, is a code.
Equality of two such configurations is \(\mathrm{EqCode}\), and \(S\) can speak about it directly.
That case is the setting of Section~\ref{sec:undecidable}.
It does not restore bit-identity as an \(S\)-predicate of nameless reals, and it does not turn Theorem~\ref{thm:poincare} into a self-certifying sentence.

\begin{definition}[Proof predicate]
\label{def:prov}
Let \(\prf_{S}(p,\gn{\varphi})\) be a primitive recursive proof relation coming from a primitive recursive enumeration of the axioms of \(S\), and set
\[
\Prov_{S}(\gn{\varphi})\ :\Longleftrightarrow\ \exists p\ \prf_{S}(p,\gn{\varphi}),
\qquad
\Con_{S}\ :\Longleftrightarrow\ \neg\Prov_{S}(\gn{0=1}).
\]
The box of this note is this predicate:
\begin{equation}
\label{eq:box}
\Box\varphi\ :\Longleftrightarrow\ \Prov_{S}(\gn{\varphi}).
\end{equation}
The Hilbert--Bernays--L\"ob derivability conditions are assumed for this choice \citep{hilbertbernays1939,lob1955,boolos1993}: for all sentences \(\varphi,\psi\),
\begin{align}
\text{(D1)}\quad & \text{if } S\vdash\varphi \text{ then } S\vdash\Prov_{S}(\gn{\varphi}), \nonumber\\
\text{(D2)}\quad & S\vdash\Prov_{S}(\gn{\varphi\to\psi})\to\bigl(\Prov_{S}(\gn{\varphi})\to\Prov_{S}(\gn{\psi})\bigr), \label{eq:hbl}\\
\text{(D3)}\quad & S\vdash\Prov_{S}(\gn{\varphi})\to\Prov_{S}\bigl(\gn{\Prov_{S}(\gn{\varphi})}\bigr). \nonumber
\end{align}
\end{definition}

Condition (D1) is a rule about \(S\), not an implication \(\varphi\to\Box\varphi\) proved inside \(S\).
Conditions (D2) and (D3) are single theorems of \(S\) for each pair of sentences, in fact theorems of \(\PA\) for the standard proof predicate.

\begin{remark}[The predicate is intensional]
\label{rem:feferman}
The second incompleteness theorem applies to a proof predicate satisfying \eqref{eq:hbl}.
It is not a theorem about every arithmetical formula that defines the same set of theorems in the standard model.
There are numerations of the axioms of a consistent \(S\) for which the associated consistency statement is provable in \(S\) \citep{feferman1960}.
Any reading that changes the proof predicate has changed the claim.
In this note \(\Prov_{S}\) always means a standard predicate as in Definition~\ref{def:prov}.
\end{remark}

\section{Modal packaging}
\label{sec:modal}

G\"odel--L\"ob logic \(\GL\) is the propositional modal logic of \(\Box\) under \eqref{eq:box} \citep{solovay1976,boolos1993}.
Modal formulas are built from propositional variables by Boolean connectives and \(\Box\).
The axioms of \(\GL\) are all propositional tautologies, the distribution axiom
\[
\Box(A\to B)\to(\Box A\to\Box B),
\]
and L\"ob's axiom
\begin{equation}
\label{eq:lob-axiom}
\Box(\Box A\to A)\to\Box A.
\end{equation}
The rules are modus ponens and necessitation: from a theorem \(A\), infer \(\Box A\).
Necessitation is a rule of proof.
It does not add the theorem \(A\to\Box A\).
The transitivity scheme \(\Box A\to\Box\Box A\) is derivable from \eqref{eq:lob-axiom}; it is not an extra assumption.

An \emph{arithmetical realization} is a map \(*\) from propositional variables to sentences of \(S\), extended by the Boolean clauses and by
\[
(\Box A)^{*}\ =\ \Prov_{S}(\gn{A^{*}}).
\]

\begin{theorem}[Solovay completeness]
\label{thm:solovay}
Let \(S\) be a consistent recursively axiomatized extension of \(\PA\), and let \(\Prov_{S}\) satisfy \eqref{eq:hbl}.
A modal sentence \(A\) is a theorem of \(\GL\) if and only if \(S\vdash A^{*}\) for every arithmetical realization \(*\).
\end{theorem}

\begin{proof}[Citation]
This is Solovay's arithmetical completeness theorem \citep{solovay1976}; a full development is \citet[Chapters~1--5]{boolos1993}.
The note uses it as a completed theorem.
It does not reprove the Solovay embedding.
\end{proof}

The scheme one would need in order to pass from ``\(S\) proves that it returns'' to ``it returns'' is reflection, \(\Box A\to A\).
That scheme is not free in \(\GL\).

\begin{proposition}[Reflection is not a theorem of \(\GL\)]
\label{prop:no-reflection}
\(\GL\nvdash \Box p\to p\).
On the Kripke frame consisting of a single world with no accessibility arrow, \(\Box p\to p\) is not valid.
Over \(S\), the arithmetical instance \(\Prov_{S}(\gn{0=1})\to 0=1\) is equivalent to \(\Con_{S}\).
\end{proposition}

\begin{proof}
Finite transitive irreflexive frames validate \(\GL\) \citep{boolos1993}.
In a one-point frame with empty accessibility relation, \(\Box p\) is vacuously true at the only world, because that world sees no world at which \(p\) fails.
Interpreting \(p\) as false at that world makes \(\Box p\to p\) false there.
Hence the scheme is not \(\GL\)-valid, and not a \(\GL\)-theorem.

For the arithmetical instance, \(\PA\) proves \(\neg(0=1)\).
Propositional logic then yields equivalence, inside \(S\), of \(\neg\Prov_{S}(\gn{0=1})\) and \(\Prov_{S}(\gn{0=1})\to 0=1\).
The left-hand side is \(\Con_{S}\).
\end{proof}

By Theorem~\ref{thm:solovay} and the second incompleteness theorem (Theorem~\ref{thm:g2} below), \(S\) does not prove every arithmetical realization of \(\Box p\to p\).
The one-point frame shows the modal failure before any appeal to arithmetic.
The two failures match.

L\"ob's axiom is the relation \(\GL\) does prove.
Read as a statement about a sentence \(R\), it says:
\begin{equation}
\label{eq:lob-reading}
\Box(\Box R\to R)\ \to\ \Box R.
\end{equation}
The hypothesis is that the reflection instance for \(R\) is itself provable in \(S\).
The conclusion is that \(R\) is provable in \(S\).
Truth of \(R\) is a separate claim.
Passing from \(\Box R\) to \(R\) uses reflection, which Proposition~\ref{prop:no-reflection} withholds.

There is a standard extension used for truth rather than for provability.
\(\GLS\) is \(\GL\) together with the axiom scheme \(\Box A\to A\), with necessitation restricted to theorems of \(\GL\) \citep{boolos1993}.
The restriction is required for consistency of the calculus: unrestricted necessitation applied to the reflection instance \(\Box\bot\to\bot\) yields \(\Box(\Box\bot\to\bot)\), L\"ob's axiom yields \(\Box\bot\), and reflection then yields \(\bot\).
Under a sound realization, theorems of \(\GLS\) are true in the standard model.
The reflection axioms of \(\GLS\) are not theorems of \(S\), and they are not theorems of \(\GL\).
External soundness of \(S\), when it is assumed, is an assumption about the standard model.
It is not a derivation inside \(S\), and it is not a hypothesis of Theorem~\ref{thm:poincare}.

\begin{remark}[``Necessarily returns'']
\label{rem:necessarily}
Take \(R\) to be an arithmetical sentence proposed as a coded return claim.
The narrative inference ``necessarily returns, therefore returns'' is the inference from \(\Box R\) to \(R\).
\(\GL\) does not license it.
The L\"ob condition under which \(S\) proves \(R\) is \(S\vdash \Prov_{S}(\gn{R})\to R\), which by L\"ob's theorem (Theorem~\ref{thm:lob}) already yields \(S\vdash R\).
In that case the box is not an independent witness of return.
\(S\) proves the return sentence outright, or the L\"ob hypothesis has not been met.
The finite-measure hypotheses of Theorem~\ref{thm:poincare} do not provide \(S\vdash \Box R\to R\).
\end{remark}

\section{Epistemic packaging}
\label{sec:epistemic}

The operator of this note is \(\Box\), defined by \eqref{eq:box} as provability in \(S\).
This section records what would change if the same box were reread as knowledge, and then declines that rereading.

The modal system \(\mathsf{S5}\) for an operator \(K\) has distribution, the factivity axiom \(KA\to A\), the positive-introspection axiom \(KA\to KKA\), and the negative-introspection axiom \(\neg KA\to K\neg KA\), with modus ponens and necessitation.
Factivity is the inference ``known return yields return.''
Proposition~\ref{prop:no-reflection} says that the corresponding scheme for \(\Prov_{S}\) is not a theorem of \(\GL\) and that its instance at \(0=1\) is \(\Con_{S}\).
A calculus that adds factivity for the same symbol that denotes \(\Prov_{S}\) is the wrong calculus for that symbol.
Calling the operator knowledge does not change the proof predicate.

A factive complex can be defined in the metatheory without pretending that \(S\) defines it:
\begin{equation}
\label{eq:known}
K_{S}(\varphi)\ :\Longleftrightarrow\ \Prov_{S}(\gn{\varphi})\ \text{is true and}\ \varphi\ \text{is true}.
\end{equation}
The implication \(K_{S}(\varphi)\to\varphi\) holds by the second conjunct, in the metatheory.
It is not an \(S\)-theorem scheme.
\(S\) cannot define a truth predicate for its own language \citep{tarski1936}, so \eqref{eq:known} is not a formula of \(S\).
``Knowledge of return,'' on this reading, is an external comment on a sentence.
It is not a configuration in \(X\), and it is not a proof of \(\Con_{S}\).

The note therefore does not adopt a knower, an \(\mathsf{S5}\) axiom, or a knowability principle.
Where a box appears below, it is \(\Prov_{S}\).

\section{Self-certification barrier}
\label{sec:barrier}

Self-certification is a particular arithmetical sentence, obtained when a configuration code is required to contain \(S\), or to contain a proof, and the sentence says of its own code that this happens.
The diagonal lemma supplies the sentence.
It does not by itself decide the sentence.

\begin{lemma}[Diagonal lemma]
\label{lem:diag}
For every formula \(\psi(v)\) of \(S\) with one free variable, there is a sentence \(\rho\) such that
\[
S\ \vdash\ \rho\ \leftrightarrow\ \psi(\gn{\rho}).
\]
\end{lemma}

\begin{proof}[Citation]
The usual fixed-point construction from the proof of the first incompleteness theorem \citep{godel1931,boolos1993}.
\end{proof}

Two fixed points have to be kept apart, and both have to be kept apart from \(\varepsilon\)-Return.

\begin{definition}[Three sentences]
\label{def:three}
Let \(R(v)\) be any formula of \(S\) with one free variable.
The intended reading of \(R(c)\), when one is fixed, is an arithmetization of ``the configuration coded by \(c\) recurs'' or ``the configuration coded by \(c\) encodes an \(S\)-proof of a named sentence.''
The barrier statements do not depend on further details of \(R\).
\begin{enumerate}
\item \textbf{External recurrence.}
The sentence \(E\) of set theory is Theorem~\ref{thm:poincare}, including the finite-measure hypotheses in its antecedent.
Its provability, if discussed, is \(\Prov_{\ZFC}(\gn{E})\).

\item \textbf{Henkin certificate.}
By Lemma~\ref{lem:diag} applied to \(\Prov_{S}(v)\), there is a sentence \(\sigma\) with
\begin{equation}
\label{eq:henkin}
S\ \vdash\ \sigma\ \leftrightarrow\ \Prov_{S}(\gn{\sigma}).
\end{equation}
This is the sentence that says ``\(S\) proves this sentence.''

\item \textbf{Consistency-bundled return.}
By Lemma~\ref{lem:diag} applied to \(R(v)\wedge\Con_{S}\), there is a sentence \(\rho\) with
\begin{equation}
\label{eq:bundled}
S\ \vdash\ \rho\ \leftrightarrow\ \bigl(R(\gn{\rho})\wedge\Con_{S}\bigr).
\end{equation}
This is the sentence that says ``the recurrence formula holds of my code, and \(S\) is consistent.''
\end{enumerate}
\end{definition}

Encoding \(S\) in a configuration is already covered by the second clause when the certificate is a proof code, and by the third when the sentence demands \(\Con_{S}\) itself.
A code for the axiom set of \(S\) can sit inside the number \(\gn{\rho}\).
That is ordinary G\"odel numbering.
It does not add an axiom to \(S\).

\begin{theorem}[G\"odel's second incompleteness theorem]
\label{thm:g2}
If \(S\) is consistent and \(\Prov_{S}\) satisfies \eqref{eq:hbl}, then \(S\nvdash\Con_{S}\).
\end{theorem}

\begin{proof}[Citation]
\citet{godel1931}, in the form isolated by the derivability conditions \citep{hilbertbernays1939,lob1955,boolos1993}.
The intensional caveat is Remark~\ref{rem:feferman}.
\end{proof}

\begin{theorem}[L\"ob]
\label{thm:lob}
If \(S\vdash \Prov_{S}(\gn{\varphi})\to\varphi\), then \(S\vdash\varphi\).
Equivalently, under \eqref{eq:hbl},
\[
S\ \vdash\ \Prov_{S}\bigl(\gn{\Prov_{S}(\gn{\varphi})\to\varphi}\bigr)\ \to\ \Prov_{S}(\gn{\varphi}).
\]
\end{theorem}

\begin{proof}[Citation]
\citet{lob1955}.
The modal form is \eqref{eq:lob-axiom}, via Theorem~\ref{thm:solovay}.
\end{proof}

\begin{proposition}[The Henkin certificate does not prove consistency]
\label{prop:henkin}
Assume \(S\) is consistent and \(\Prov_{S}\) is standard.
Then \(S\vdash\sigma\), the sentence \(\sigma\) is true in the standard model, \(\Con_{S}\) is true in the standard model, and
\[
S\ \nvdash\ \sigma\to\Con_{S}.
\]
\end{proposition}

\begin{proof}
From \eqref{eq:henkin}, \(S\vdash \Prov_{S}(\gn{\sigma})\to\sigma\).
Theorem~\ref{thm:lob} gives \(S\vdash\sigma\).
A standard proof predicate represents actual provability, so \(\Prov_{S}(\gn{\sigma})\) is true in \(\Nat\), and therefore \(\sigma\) is true in \(\Nat\).
Consistency of \(S\) means that \(0=1\) has no \(S\)-proof, so the standard sentence \(\Con_{S}\) is true in \(\Nat\).
The implication \(\sigma\to\Con_{S}\) is therefore true in \(\Nat\).

It is not a theorem of \(S\).
If \(S\vdash\sigma\to\Con_{S}\), then \(S\vdash\sigma\) and modus ponens would give \(S\vdash\Con_{S}\), contradicting Theorem~\ref{thm:g2}.
\end{proof}

Proposition~\ref{prop:henkin} is the precise status of ``this sentence is proved in \(S\).''
\(S\) does prove that self-certificate.
Truth of the certificate in \(\Nat\) is secured by the existence of the proof L\"ob's theorem constructs.
What fails inside \(S\) is the passage from that certificate to \(\Con_{S}\).
The implication can be true in \(\Nat\) and still unprovable in \(S\).
That is incompleteness of \(S\) at a true sentence.
It is not a contradiction between the certificate and consistency, and it is not a statement about points of \(X\).

\begin{proposition}[The bundled fixed point is unprovable]
\label{prop:bundled}
Assume \(S\) is consistent.
Let \(\rho\) be as in \eqref{eq:bundled} for an arbitrary formula \(R(v)\).
Then \(S\vdash\rho\to\Con_{S}\), and \(S\nvdash\rho\).
If \(R(v)\) is the formula \(v=v\), then \(\rho\) is equivalent over \(S\) to \(\Con_{S}\), hence true in \(\Nat\) and unprovable in \(S\).
If \(R(v)\) is the formula \(v\neq v\), then \(S\vdash\neg\rho\).
For a general recurrence formula, unprovability of \(\rho\) holds either way, and the truth-value of \(\rho\) in \(\Nat\) tracks the truth-value of \(R(\gn{\rho})\).
\end{proposition}

\begin{proof}
The left-to-right direction of \eqref{eq:bundled} is \(S\vdash\rho\to\Con_{S}\).
If \(S\vdash\rho\), modus ponens yields \(S\vdash\Con_{S}\), contradicting Theorem~\ref{thm:g2}.
Hence \(S\nvdash\rho\).

If \(S\vdash R(\gn{\rho})\leftrightarrow \top\), then \eqref{eq:bundled} gives \(S\vdash\rho\leftrightarrow\Con_{S}\).
Consistency makes \(\Con_{S}\) true and Theorem~\ref{thm:g2} makes it unprovable, so the same holds of \(\rho\).
If \(S\vdash R(\gn{\rho})\leftrightarrow \bot\), then \(S\vdash\neg\rho\).
In the general case, \(\Con_{S}\) is true, so \(\Nat\models\rho\) if and only if \(\Nat\models R(\gn{\rho})\).
The note does not evaluate \(R(\gn{\rho})\) for any particular arithmetization of a dynamical return.
\end{proof}

The slogan ``this life returns, including this proof that it returns'' collapses Definition~\ref{def:three}.
Split into the sentences actually available, it becomes one of the following, and they have different statuses.

\begin{center}
\begin{tabular}{@{}llp{6.4cm}@{}}
\toprule
Reading & Sentence & Status, when \(S\) is consistent \\
\midrule
\(\varepsilon\)-Return & \(E\) in \(\ZFC\) & Classical theorem under finite measure. Not a \(\Con_{S}\) claim. \\
Henkin certificate & \(\sigma\leftrightarrow\Box\sigma\) & Provable in \(S\). Does not prove \(\Con_{S}\). \\
Bundled return & \(\rho\leftrightarrow R(\gn{\rho})\wedge\Con_{S}\) & Unprovable in \(S\). Truth tracks \(R\), not a refutation of return. \\
Detachment & \(\Box A\to A\) & Not a theorem of \(\GL\). The instance \(A=\bot\) is \(\Con_{S}\). \\
\bottomrule
\end{tabular}
\end{center}

\begin{remark}[No exemption by coding a configuration]
\label{rem:no-exemption}
Theorem~\ref{thm:g2} has no side condition of the form ``unless a configuration code contains an index of \(S\).''
A numeral denoting such a code can appear in a sentence of \(S\).
If that sentence is \(\Con_{S}\) itself, \(S\) does not prove it.
If that sentence is \(\sigma\), Proposition~\ref{prop:henkin} applies.
If that sentence is \(\rho\), Proposition~\ref{prop:bundled} applies.
Writing a description of \(S\) into a number does not produce a derivation of \(\Con_{S}\) from the axioms of \(S\).
\end{remark}

\begin{remark}[Witnessed \(\Sigma_1\) returns are a different sentence]
\label{rem:sigma1}
A single coded orbit is not the bundled fixed point.
If \(\varphi\) is a \(\Sigma_1\) sentence and \(\Nat\models\varphi\), then \(\PA\vdash\varphi\).
In particular, if a computable configuration \(C\) actually returns to itself, a computation trace is a witness, and the \(\Sigma_1\) sentence ``\(C\) recurs'' is a theorem of \(\PA\).
Incompleteness does not hide witnessed returns.
The sentences that self-certification adds are of another kind: \(\Con_{S}\) is \(\Pi_1\), reflection at a general sentence is not a true \(\Sigma_1\) fact supplied by a return time, and \(\rho\) conjoins a recurrence formula with \(\Con_{S}\).
\(\Sigma_1\)-completeness settles the witnessed \(\Sigma_1\) sentence and does not settle \(\rho\).
\end{remark}

\begin{remark}[A stronger theory is not an internal certificate]
\label{rem:gentzen}
Theorem~\ref{thm:g2} says that \(S\) does not prove \(\Con_{S}\).
It does not say that no theory proves \(\Con_{S}\).
A stronger system can prove the consistency of a weaker one.
Whatever that stronger proof is, it is a derivation in the stronger system.
It is not a derivation in \(S\), and placing a copy of it inside a described configuration does not make it a derivation in \(S\).
The barrier in this section is self-certification: proving \(\Con_{S}\) from within \(S\), or detaching \(\Box\) inside \(\GL\).
\end{remark}

Nothing in Propositions~\ref{prop:henkin} and~\ref{prop:bundled} mentions \(\mu\), \(T\), or a neighborhood.
They remain true if the formula \(R(v)\) arithmetizes a fact unrelated to dynamics.
They therefore cannot cancel a measure calculation in \(\ZFC\), and they cannot supply the finite-measure hypothesis that Theorem~\ref{thm:poincare} requires.

\section{Undecidability of recurrence queries}
\label{sec:undecidable}

The question ``does configuration \(C\) recur?'' is a decision problem about a named class of systems.
Its undecidability, when it holds, is a fact about that decision problem.

\begin{definition}
\label{def:config}
A \emph{Turing configuration} is a triple \((q,t,h)\) where \(q\) is a state of a Turing machine, \(t\) is a tape inscription of finite support, and \(h\in\ZZ\) is a head position.
\(\mathrm{Step}_{M}\) is the usual one-step function on configurations of a machine \(M\).
A configuration \(C\) \emph{recurs} under \(M\) when there exists \(n\geq 1\) with \(\mathrm{Step}_{M}^{n}(C)=C\).
\end{definition}

Equality here is equality of finitely presented configurations.
In the sense of Definition~\ref{def:codes}, it is \(\mathrm{EqCode}\), not neighborhood membership in a continuum.

\begin{proposition}[Recurrence queries can encode halting]
\label{prop:undecidable}
The set of pairs \((M,C)\) such that \(C\) recurs under \(M\) is not recursive.
\end{proposition}

\begin{proof}[Standard reduction]
The halting set \(\{(N,w) \mid N \text{ halts on input } w\}\) is not recursive \citep{turing1936}.
From a pair \((N,w)\), construct a machine \(M_{N,w}\) whose finite control contains a copy of the transition table of \(N\) and a procedure that writes the finite word \(w\).
The map \((N,w)\mapsto M_{N,w}\) is computable.
Let \(q_0\) be the start state of \(M_{N,w}\) and let \(C_0\) be the configuration \((q_0,\text{blank tape},0)\).
In the transition table, \(q_0\) occurs only as the source of the first step and as the target of the cleanup step described below.

From \(C_0\) the first transition enters a work state distinct from \(q_0\).
The machine writes \(w\), a separator, and a binary counter, then simulates \(N\) on \(w\), incrementing the counter once per simulated step.
Every configuration in this phase is distinct from \(C_0\): the state is not \(q_0\).
If the simulation halts, the machine enters a cleanup state \(q_{\mathrm{clean}}\neq q_0\) and sweeps the finite marked region, erasing symbols until the tape is blank and the head is on cell \(0\), remaining in \(q_{\mathrm{clean}}\) throughout the sweep.
One further transition sends that blank-tape configuration to \(C_0\).
A single Turing step is not asked to erase an arbitrary tape; the sweep does the erasing, and only the last step enters \(q_0\).

If \(N\) halts on \(w\), the run from \(C_0\) reaches that last step, so \(C_0\) recurs.
If \(N\) does not halt on \(w\), the cleanup state is never entered, the counter is unbounded, and no later configuration has a blank tape in state \(q_0\).
Thus \(C_0\) does not recur.
A decision procedure for recurrence of pairs \((M,C)\) would decide the halting set by computing \((M_{N,w},C_0)\).
\end{proof}

The same style of encoding is what makes various long-term questions about dynamical systems that contain a universal computation undecidable \citep{moore1990}.
Proposition~\ref{prop:undecidable} is the special case in which the long-term question is exact return of a finitely presented configuration.
The short proof above is the direct reduction; it does not depend on the extra structure of generalized shifts.

\begin{remark}[What the reduction produces]
\label{rem:both}
If \(N\) halts, the constructed \(C_0\) recurs.
If \(N\) does not halt, that \(C_0\) does not recur.
Both sorts of instance exist: there are machines that halt and machines that do not.
Undecidability is the nonexistence of a total algorithm that sorts all instances.
It is not an assertion that recurrent configurations do not exist, and it is not an assertion that they do.
\end{remark}

\begin{proposition}[Finitely presented systems are decidable]
\label{prop:finite}
Let \(X\) be a finite set and let \(T\colon X\to X\) be given explicitly.
For every \(x\in X\) there exist \(k\geq 0\) and \(n\geq 1\) with \(k+n\leq |X|\) and \(T^{k+n}(x)=T^{k}(x)\).
The point \(x\) itself recurs if and only if \(T^{n}(x)=x\) for some \(n\) with \(1\leq n\leq |X|\).
Both questions are decidable by running the orbit for at most \(|X|\) steps.
\end{proposition}

\begin{proof}
The \(|X|+1\) configurations \(x,Tx,\ldots,T^{|X|}x\) are not pairwise distinct, so \(T^{a}(x)=T^{b}(x)\) for some \(0\leq a<b\leq |X|\).
The first repetition closes a cycle.
Running \(|X|\) steps finds it.
\end{proof}

Proposition~\ref{prop:finite} is the reason Proposition~\ref{prop:undecidable} needs an unbounded tape.
The counter in the reduction takes infinitely many distinct values along a non-halting simulation.
The counting measure on the set of all finite-support configurations is an infinite measure.
Theorem~\ref{thm:poincare} does not apply to that measure, so a non-recurrent \(C_0\) is not a counterexample to \(\varepsilon\)-Return.
A system whose configuration set is finite falls under Proposition~\ref{prop:finite}: every orbit eventually cycles, and the cycle is found by inspection.
Exact eventual periodicity in a finite set is still \(\mathrm{EqCode}\) on that set.
It is not a proof of \(\Con_{S}\), and it is not \(\varepsilon\)-Return on a continuum.

\section{What is not claimed}
\label{sec:not}

The positive content of the note is the separation recorded above.
The following consequences do not follow from it.

\begin{enumerate}
\item \textbf{\(\varepsilon\)-Return stands under its own hypotheses.}
Theorem~\ref{thm:poincare} is a theorem of ordinary mathematics from finite measure and invariance.
Its formalization is a theorem of \(\ZFC\) under those hypotheses written into the antecedent.
Consistency or incompleteness of \(S\) is not a hypothesis of the proof in Section~\ref{sec:setup}, and it does not appear in the conclusion.
Incompleteness of \(S\) does not entail that orbits fail to return.

\item \textbf{Incompleteness does not support the recurrence theorem.}
Theorem~\ref{thm:g2} says \(S\nvdash\Con_{S}\).
It does not enlarge \(\mu\), it does not force \(\mu(X)<\infty\), and it does not move a point into a neighborhood of itself.
There is no inference from \(S\nvdash\Con_{S}\) to \(\varepsilon\)-Return, nor from \(S\nvdash\Con_{S}\) to exact identity of states.
A failure of self-certification is not evidence that a phase point recurs.

\item \textbf{The bundled sentence is unprovable, and its truth tracks \(R\).}
Proposition~\ref{prop:bundled} says that consistent \(S\) does not prove \(\rho\).
If the recurrence formula \(R(\gn{\rho})\) is true in \(\Nat\), then \(\rho\) is true and unprovable.
If it is false, then \(\rho\) is false.
The note supplies no dynamical calculation that would decide \(R(\gn{\rho})\).
Unprovability of the bundle is not a disproof of recurrence.

\item \textbf{Witnessed returns remain provable when they are \(\Sigma_1\) and true.}
Remark~\ref{rem:sigma1} is part of the claim.
A trace of a return of a coded configuration is a proof in \(\PA\) of that \(\Sigma_1\) sentence.
The self-certification barrier concerns certificates that ask \(S\) to prove \(\Con_{S}\), or ask \(\GL\) to detach \(\Box A\to A\), or ask a sentence to include \(\Con_{S}\) in its own fixed point.

\item \textbf{\(\Box\) is provability in \(S\).}
It is not an \(\mathsf{S5}\) knowledge operator and not a metaphysical necessity.
``Necessarily returns, therefore returns'' is \(\Box R\to R\), which Proposition~\ref{prop:no-reflection} excludes from \(\GL\).
The admissible substitute is L\"ob's condition \(S\vdash\Box R\to R\), which yields \(S\vdash R\) rather than a new passage from a box to a fact about \(X\).

\item \textbf{\(\varepsilon\)-Return is not the demand to will an identical life.}
The published ethical sentence asks whether one can will this life in the same order, again and innumerable times \citep[\S 341]{nietzsche1882}.
\(\varepsilon\)-Return says that, under named measure hypotheses, almost every state returns to every neighborhood of itself.
A neighborhood of a state is membership in a positive-radius set.
Proposition~\ref{prop:rotation} meets the measure hypotheses at every point and meets exact identity at none.
Uniqueness of one biographical sequence is not a conclusion of Theorem~\ref{thm:poincare}.
The ethical sentence is not a theorem of \(\ZFC\) derived here, and it is not a theorem of \(\GL\).

\item \textbf{G\"odel 1949 is a different construction.}
The rotating universes of \citet{godel1949} are exact solutions of Einstein's field equations in which some worldlines are closed.
That is a geometric statement under those field equations and those metrics.
It is not \(\varepsilon\)-Return, it is not \(\Con_{S}\), and it is not Proposition~\ref{prop:undecidable}.

\item \textbf{Undecidability of a query is not a theorem about a measure.}
Proposition~\ref{prop:undecidable} says that no algorithm classifies all machine-configuration pairs.
Remark~\ref{rem:both} records that the reduction builds recurrent configurations and non-recurrent ones.
Proposition~\ref{prop:finite} records that a fully presented finite system is decidable by inspection and eventually periodic.
None of these sentences is the negation of Theorem~\ref{thm:poincare}, and none of them is that theorem.
\end{enumerate}

The three readings in the introduction can now be stated in one place, with the systems named.
\(\varepsilon\)-Return is Theorem~\ref{thm:poincare} and Corollary~\ref{cor:epsilon}, proved in analysis from finite measure, and formalizable in \(\ZFC\).
What \(S\) can say about a continuous state is \(\mathrm{EqCode}\) or \(\mathrm{Near}\) for a chosen code, as in Remark~\ref{rem:boundary}.
What \(S\) can say about a certificate that includes \(S\) is Proposition~\ref{prop:henkin} for the Henkin sentence and Proposition~\ref{prop:bundled} for the sentence that conjoins \(\Con_{S}\); the passage from \(\Box R\) to \(R\) is Remark~\ref{rem:necessarily}.
The decision problem for exact recurrence of Turing configurations is Proposition~\ref{prop:undecidable}.
Each sentence keeps the hypotheses of its own proof.
None of them is a booster for another.

\bibliographystyle{plainnat}
\bibliography{refs}

\end{document}
